\documentclass[11pt]{article}
\usepackage[margin=1in]{geometry}
\usepackage{amsmath,amssymb,amsthm,mathtools}
\usepackage{graphicx,xcolor,enumitem,microtype,hyperref}
\hypersetup{colorlinks=true,linkcolor=blue!55!black,citecolor=blue!55!black,urlcolor=blue!55!black}
\newtheorem{theorem}{Theorem}
\newtheorem{lemma}{Lemma}
\newtheorem{corollary}{Corollary}
\newcommand{\A}{\mathcal A}
\newcommand{\C}{\mathbb C}
\newcommand{\dbar}{\overline\partial}
\newcommand{\bd}{\partial}
\newcommand{\cl}[1]{\overline{#1}}
\newcommand{\K}{K^2_{(0,q)}}
\newcommand{\Ho}{H^q}
\title{\vspace{-2em}Continuous-Symbol Hankel Compactness\\on Product Domains}
\author{Partial solution packet for arXiv:2011.02656}
\date{17 August 2026}

\begin{document}
\maketitle
\begin{center}
\fcolorbox{orange!75!black}{orange!4}{\parbox{0.91\textwidth}{
\textbf{Status: candidate substantial partial theorem, likely valid.}
Remark~2 of the source has an affirmative answer, at every form level, on
products of two bounded strictly convex domains.  It also has an affirmative
$q=0$ answer on arbitrary finite products of bounded convex domains.  The
general convex-domain question remains open.
}}
\end{center}

\section{Source question and result}

Let $\Omega\subset\C^n$ be a bounded convex domain.  For
$0\leq q\leq n-1$, write $K^2_{(0,q)}(\Omega)$ for the square-integrable
$\dbar$-closed $(0,q)$-forms, $P_q$ for the orthogonal projection onto that
space, and
\[
  H_\varphi^q f=(I-P_q)(\varphi f).
\]
Celik, Sahutoglu, and Straube prove that if
$\varphi\in C^1(\cl\Omega)$ is holomorphic along all analytic varieties in
$\bd\Omega$ of dimension at least $q+1$, then $H_\varphi^q$ is compact
\cite[Theorem~1]{CSS}.  Their Remark~2 asks whether $C^1(\cl\Omega)$ can be
replaced by $C(\cl\Omega)$.

\begin{figure}[ht]
\centering
\includegraphics[width=0.98\textwidth]{figures/open_problem_crop.png}
\caption{The source theorem and open continuous-symbol question
on source page~2.  The displayed page was reconstructed from the run's
ingested official arXiv source.}
\end{figure}

Our main result settles that question on a product class with uncountably
many side varieties.

\begin{theorem}[Two strictly convex factors]\label{thm:main}
Let $D\subset\C^a$ and $G\subset\C^b$ be bounded strictly convex domains,
$\Omega=D\times G$, and $0\leq q\leq a+b-1$.  Suppose
$\varphi\in C(\cl\Omega)$ is holomorphic along every analytic variety in
$\bd\Omega$ of dimension at least $q+1$.  Then
\[
       H_\varphi^q:K^2_{(0,q)}(\Omega)\longrightarrow
       L^2_{(0,q)}(\Omega)
\]
is compact.
\end{theorem}

For $q=0$, strict convexity of the factors is unnecessary and any finite
number of factors is allowed; see Corollary~\ref{cor:finite}.

\section{Two approximation principles}

We repeatedly use the elementary operator bound
\begin{equation}\label{eq:norm}
 \|H_\varphi^q-H_\psi^q\|
 =\|(I-P_q)M_{\varphi-\psi}|_{K^2_{(0,q)}}\|
 \leq\|\varphi-\psi\|_\infty .
\end{equation}
Thus a uniform limit of symbols with compact Hankel operators again has a
compact Hankel operator.

For a bounded domain $U$, put
\[
       A(U)=C(\cl U)\cap\mathcal O(U)
\]
with the uniform norm.  The following boundary-cross lemma is the central
device.

\begin{lemma}[Boundary cross]\label{lem:cross}
Let $D\subset\C^a$ and $G\subset\C^b$ be bounded domains and let
$\varphi\in C(\cl D\times\cl G)$.  Assume
\[
\begin{aligned}
 &w\longmapsto\varphi(\xi,w)\in A(G)
       &&(\xi\in\bd D),\\
 &z\longmapsto\varphi(z,\eta)\in A(D)
       &&(\eta\in\bd G).
\end{aligned}
\]
Then there is a unique $h\in A(D\times G)$ satisfying
$h=\varphi$ on $\bd(D\times G)$.
\end{lemma}

\begin{proof}
Let $B_D\subset C(\bd D)$ be the boundary traces of $A(D)$.  Restriction
$T_D:A(D)\to B_D$ is an isometry by the maximum principle.  Hence $B_D$ is
closed and $R_D=T_D^{-1}$ is an isometry.

Define $F:\cl G\to C(\bd D)$ by
$F(w)(\xi)=\varphi(\xi,w)$.  Uniform continuity of $\varphi$ makes $F$
continuous.  It is Banach-valued holomorphic on $G$: for every finite
measure $\mu$ on $\bd D$, Morera's theorem and Fubini show that
$w\mapsto\int_{\bd D}\varphi(\xi,w)\,d\mu(\xi)$ is holomorphic.

If $\Lambda\in C(\bd D)^*$ annihilates $B_D$, then
$\Lambda(F(\cdot))\in A(G)$.  For $\eta\in\bd G$ the second slice assumption
gives $F(\eta)\in B_D$, so this scalar function vanishes on $\bd G$.
The maximum principle gives $\Lambda(F(w))=0$ on $G$.  Hahn--Banach now
implies $F(w)\in B_D$ for every $w\in\cl G$.

Set $H(w)=R_DF(w)\in A(D)$ and $h(z,w)=H(w)(z)$.  The map $H$ is continuous
on $\cl G$ and Banach-holomorphic on $G$.  Thus $h$ is separately
holomorphic, hence jointly holomorphic by Hartogs' theorem, and is continuous
on the closed product.  Its trace agrees with $\varphi$ on
$\bd D\times\cl G$ by construction.  On $\cl D\times\bd G$, the two
$A(D)$ functions have the same boundary trace and hence agree.  Uniqueness
follows from the maximum principle.
\end{proof}

We also need a one-sided version at the level of approximation.

\begin{lemma}[One-sided smooth approximation]\label{lem:one}
Let $D\subset\C^a$ be bounded and convex, $G\subset\C^b$ bounded, and
$h\in C(\cl D\times\cl G)$ be holomorphic in $z\in D$ for each
$w\in\cl G$.  Then $h$ is a uniform limit of functions in
$C^1(\cl D\times\cl G)$ that are holomorphic in $z$.

If $g\in C(\cl D\times\cl G)$ vanishes on
$\cl D\times\bd G$, then it is a uniform limit of $C^1$ functions that
vanish on a neighborhood of that side.
\end{lemma}

\begin{proof}
Regard $h$ as a continuous map $\cl G\to A(D)$.  Polynomials are dense in
$A(D)$: dilate toward an interior point of the convex domain and then apply
Oka--Weil.  The compact image $h(\cl G)$ therefore has a finite cover by
small balls centered at polynomials $p_j$.  A smooth partition of unity
$\chi_j$ on a neighborhood of $\cl G$, subordinate to the inverse-image
cover, gives
\[
                         h_\varepsilon(z,w)
                  =\sum_j\chi_j(w)p_j(z),
\]
with $\|h-h_\varepsilon\|_\infty<\varepsilon$.

For the second assertion, uniform continuity and vanishing on the indicated
side allow multiplication by a continuous cutoff in $w$ that is zero near
$\bd G$ with arbitrarily small uniform error.  Approximate the resulting
continuous function uniformly on the compact product by smooth functions
(real-coordinate polynomials suffice), and multiply by a slightly larger
smooth cutoff that is identically one on the first cutoff's support.  The
approximants remain zero near $\cl D\times\bd G$.
\end{proof}

The same cutoff argument shows that a continuous function vanishing on the
full boundary of a bounded domain is uniformly approximable by $C^1$
functions supported away from the boundary.  Also, every continuous function
on a compact subset of Euclidean space is uniformly approximable by smooth
functions.

\section{Boundary geometry of a strict product}

\begin{lemma}\label{lem:geometry}
Every positive-dimensional analytic boundary variety of
$D\times G$, with $D$ and $G$ strictly convex, is contained in a side
\[
             \{\xi\}\times\cl G\quad(\xi\in\bd D),
       \qquad\text{or}\qquad
             \cl D\times\{\eta\}\quad(\eta\in\bd G).
\]
Consequently its dimension is at most $b$ in the first case and at most
$a$ in the second.
\end{lemma}

\begin{proof}
Consider a connected irreducible local parameterization
$v=(v_D,v_G)$ of the variety.  Both coordinate maps take values in the
corresponding closed domains.  If $v_D$ reaches a point of $\bd D$, choose a
real supporting hyperplane there.  Its complex-linear defining functional,
composed with $v_D$, has real part bounded above and attaining its maximum at
an interior parameter point.  The maximum principle places the entire image
of $v_D$ in that supporting face.  Strict convexity makes the face a
singleton, so $v_D$ is constant at a boundary point.  If $v_D$ never reaches
$\bd D$, then $v_D$ lies in $D$; because $v$ lies in $\bd(D\times G)$,
$v_G$ lies in $\bd G$ everywhere, and the symmetric argument makes $v_G$
constant.  Applying this componentwise proves the claim and the dimension
bounds.
\end{proof}

This argument is also consistent with the source's observation that boundary
varieties of convex domains are affine intersections \cite{CSS,CSS20}.

\section{Proof of the main theorem}

\begin{proof}[Proof of Theorem~\ref{thm:main}]
There are three cases.

\smallskip
\noindent\emph{Case 1: $a,b\geq q+1$.}
For every $\xi\in\bd D$, the side $\{\xi\}\times G$ is a $b$-dimensional
boundary variety, so $\varphi(\xi,\cdot)\in A(G)$.  Similarly
$\varphi(\cdot,\eta)\in A(D)$ for every $\eta\in\bd G$.
Lemma~\ref{lem:cross} supplies $h\in A(D\times G)$ agreeing with $\varphi$
on the full boundary.  Put $g=\varphi-h$.  Then $H_h^q=0$, because
multiplication by a holomorphic function preserves $\dbar$-closedness.
Approximate $g$ uniformly by $g_\nu\in C^1(\cl\Omega)$ supported away from
$\bd\Omega$.  Each $g_\nu$ is zero along every boundary variety, so the
source theorem makes $H_{g_\nu}^q$ compact.  Equation~\eqref{eq:norm} gives
compactness of $H_g^q$ and hence of $H_\varphi^q$.

\smallskip
\noindent\emph{Case 2: exactly one of $a,b$ is at least $q+1$.}
Assume $a\geq q+1>b$; the other case is symmetric.  Lemma~\ref{lem:geometry}
shows that every boundary variety of dimension at least $q+1$ lies in
$D\times\{\eta\}$ for some $\eta\in\bd G$.  Hence
\[
             F(\eta)=\varphi(\cdot,\eta)\in A(D),
             \qquad \eta\in\bd G,
\]
and $F:\bd G\to A(D)$ is continuous.

Extend $F$ continuously to $\cl G$ with values in $A(D)$.  For completeness,
one explicit extension is obtained by choosing $w_0\in G$ and an
$F_0\in A(D)$, then interpolating linearly in $A(D)$ from $F_0$ to the
boundary value along each ray from $w_0$; convexity of $G$ makes the radial
gauge continuous.  Denote the associated scalar function by $h(z,w)$.
It is continuous on the closed product, holomorphic in $z$, and agrees with
$\varphi$ on $\cl D\times\bd G$.

By Lemma~\ref{lem:one}, choose $C^1$ functions $h_\nu$, holomorphic in $z$,
with $h_\nu\to h$ uniformly.  The remainder $g=\varphi-h$ vanishes on
$\cl D\times\bd G$; choose $C^1$ functions $g_\nu\to g$ uniformly that
vanish near that side.  Both $h_\nu$ and $g_\nu$ are holomorphic along every
boundary variety of dimension at least $q+1$: the former are holomorphic in
the only possible varying coordinate, and the latter are zero there.
The source theorem makes $H_{h_\nu}^q$ and $H_{g_\nu}^q$ compact.
Equation~\eqref{eq:norm}, applied to their sums, proves compactness of
$H_\varphi^q$.

\smallskip
\noindent\emph{Case 3: $a,b\leq q$.}
Lemma~\ref{lem:geometry} says there are no boundary varieties of dimension
at least $q+1$.  Approximate $\varphi$ uniformly by arbitrary $C^1$
functions.  The source hypothesis is vacuous for every approximant, so the
source theorem and \eqref{eq:norm} finish the proof.
\end{proof}

\section{A finite-product corollary}

\begin{corollary}\label{cor:finite}
Let $\Omega=\Omega_1\times\cdots\times\Omega_m$, $m\geq2$, where every
$\Omega_j$ is a bounded convex domain of positive dimension.  If
$\varphi\in C(\cl\Omega)$ is holomorphic along every positive-dimensional
analytic variety in $\bd\Omega$, then the scalar Bergman-space Hankel
operator $H_\varphi^0$ is compact.
\end{corollary}

\begin{proof}
Each side obtained by fixing one coordinate at a boundary point is a
positive-dimensional boundary variety, so the restriction of $\varphi$ to
it is holomorphic.  Group $\Omega_1$ as $D$ and the product of the remaining
factors as $G$.  The two slice conditions of Lemma~\ref{lem:cross} hold:
if a point of $\bd G$ has several boundary coordinates, any one of them
places the entire $D$-slice in a side face.  Thus
$\varphi=h+g$ with $h\in A(\Omega)$ and $g=0$ on $\bd\Omega$.
Now $H_h^0=0$, and smooth interior-supported approximation of $g$, followed
by the source theorem and \eqref{eq:norm}, proves the claim.
\end{proof}

\section{Checks, novelty, and limitation}

\paragraph{Logical checks.}
The form-level proof never asserts that multiplication by an
interior-supported function is compact on $K^2_{(0,q)}$; that assertion would
be false in general for $q>0$.  Instead every smooth approximant is fed into
the source theorem.  The only closure step is the elementary operator-norm
estimate \eqref{eq:norm}.  The boundary-cross argument uses no boundary
regularity.

\paragraph{Bounded novelty search.}
The run indexes and bounded title/keyword searches found the earlier theorem
for finitely many disjoint boundary varieties \cite{CSS20}, results on
products of Hankel operators on the polydisc, and continuous-symbol results
for convex Reinhardt domains under additional hypotheses.  No theorem
matching Theorem~\ref{thm:main} or Corollary~\ref{cor:finite} was found.
The present product has uncountably many side varieties, so it is not covered
by the finite-variety theorem.  Novelty is plausible, not certified.

\paragraph{Limitation.}
The packet does not solve Remark~2 for an arbitrary bounded convex domain.
Without strict product geometry, affine boundary faces can be nested and
accumulate, and the unresolved step is a compatible smooth approximation
across that entire face system.

\begin{thebibliography}{9}
\bibitem{CSS}
M.~Celik, S.~Sahutoglu, and E.~J. Straube,
\emph{A sufficient condition for compactness of Hankel operators},
arXiv:2011.02656 (2020),
\url{https://arxiv.org/abs/2011.02656}.

\bibitem{CSS20}
M.~Celik, S.~Sahutoglu, and E.~J. Straube,
\emph{Compactness of Hankel operators with continuous symbols on convex
domains}, Houston J. Math. \textbf{46} (2020), no.~4, 1005--1016;
arXiv:2005.14323.
\end{thebibliography}

\end{document}
